% Generated draft: requires mathematical and source review.
\documentclass[11pt]{article}
\usepackage[margin=25mm]{geometry}
\usepackage{fontspec}
\setmainfont{Latin Modern Roman}
\usepackage{amsmath,amssymb,mathtools,unicode-math}
\setmathfont{Latin Modern Math}
\usepackage{enumitem,needspace,xurl,hyperref}
\setlength{\parindent}{0pt}
\setlength{\parskip}{5pt}
\setlength{\emergencystretch}{3em}
\newcommand{\N}{\mathbb{N}}
\newcommand{\Z}{\mathbb{Z}}
\newcommand{\Q}{\mathbb{Q}}
\newcommand{\R}{\mathbb{R}}
\newcommand{\C}{\mathbb{C}}
\newcommand{\F}{\mathbb{F}}
\newcommand{\A}{\mathbb{A}}
\newcommand{\PP}{\mathbb{P}}
\newcommand{\E}{\mathbb{E}}
\newcommand{\eps}{\varepsilon}
\newcommand{\dd}{\,\mathrm{d}}
\begin{document}
\section*{UnsolvedMath ID: 30005130}
\textbf{Source problem code:} catalogue entry
\subsection{Definitions and mathematical statement}
The supplied record for permanent UnsolvedMath ID \(30005130\), entitled ``Comparison Maps in p-Adic Cohomology,'' has statement status \texttt{unrecoverable}. It does not determine a complete mathematical conjecture.

The first fragment is introduced by ``conjecture: (1)'' and contains the following chain, transcribed without resolving the damaged cohomology notation:
\[
\mathsf{H\ell n}(Y)=\mathsf{H\ell n}(Y(\epsilon))
\longrightarrow \mathsf{H n}(\pi^{-1}(Y(\epsilon)))
\longrightarrow \mathsf{H\ell n}(Z).
\]
Here \(\mathsf{H\ell n}\) and \(\mathsf{H n}\) represent the printed character sequences, not defined cohomology functors. The record does not define \(Y\), \(Y(\epsilon)\), \(\epsilon\), \(\pi\), \(\ell\), or \(n\), specify the cohomology theory or coefficients, define the arrows, or state the property conjectured of this chain. No domains or quantifiers for these symbols can be recovered from the supplied text.

The passage then switches to ``Theorem 3.'' Its surviving hypotheses describe \(Z\) as a smooth proper variety over a finite extension of a field printed as \(\texttt{Fp ((p}\flat\texttt{))}\), with residue field \(k\), and having a semistable formal model also printed as \(Z\), with log special fiber \(Z_0\). The field notation is unresolved, and the repeated use of \(Z\) does not distinguish the variety from its formal model. This theorem fragment ends with ``Then we have a canonical quasi-isomorphism,'' without giving its source, target, or formula.

The supplied verification describes this passage as a damaged splice between a conjectural comparison-map fragment and a proved theorem. It does not identify a recoverable problem. The two fragments therefore cannot be combined into a precise conjecture, and the missing quasi-isomorphism cannot be supplied from this record.
\subsection{Short English statement}
The record preserves an incomplete cohomological comparison chain followed by an incomplete theorem about a smooth proper variety with semistable reduction. Essential definitions and the intended conclusion are missing. Its supplied ``open'' label records unresolved extraction and does not establish that an identifiable mathematical conjecture is open.
\subsection{Sources}
\begin{itemize}
\item Supplied local UnsolvedMath registry record, permanent ID \(30005130\), source problem code \texttt{catalogue entry}, title ``Comparison Maps in p-Adic Cohomology,'' in the Oberwolfach Reports Open Problems collection. No individual proposer is supplied; the recorded proposed year is \(2022\).
\item Supplied report citation: \emph{Algebraic K-Theory (2023)}, Oberwolfach Reports, DOI \texttt{10.4171/owr/2022/24}, pp.\ 1313--1376. Source URL and supplied literature-source link labelled ``Oberwolfach report'': \url{https://doi.org/10.4171/owr/2022/24}. This citation is supplied provenance; the report has not been independently consulted for this draft.
\item Supplied literature assessment, dated exactly \texttt{2026-08-22}: the extracted record contains only the label conjecture without an identifiable mathematical statement, and the unresolved record is therefore normalized to ``open.'' This is supplied dated curation, not an independently verified literature assessment. The background explicitly requests independent verification of time-sensitive status and cited claims.
\item Supplied statement verification: the source passage is described as a damaged splice between a conjectural comparison-map fragment and a proved theorem; no separate clean statement is promoted, and the original record remains visible for provenance. The supplied statement status is \texttt{unrecoverable}.
\end{itemize}
\end{document}
